A força radial no eixo de saída do redutor é dada pela equação:

\begin{equation}
  F_{r_t} = \frac{G_{\text{tambor}} + G_{\text{cabo}}}{2} + T \cdot 10
  \label{eq:carga_radial_no_eixo_de_saida_do_redutor}
\end{equation}

Onde:

\begin{itemize}[leftmargin=1.25cm]
  \item $F_{r_t}$: força radial no eixo de saída do redutor em $\si{\newton}$.
  \item $G_{\text{tambor}}$: peso do tambor em $\si{\newton}$ \eqref{eq:peso_tambor}.
  \item $G_{\text{cabo}}$: peso do cabo de aço em $\si{\newton}$ \eqref{eq:peso_cabo}.
  \item $T$: força de tração no cabo em $\si{\deca\newton}$ \eqref{eq:forca_de_tracao_no_cabo}.
\end{itemize}

\vspace{1em}

Substituindo os valores na equação \eqref{eq:carga_radial_no_eixo_de_saida_do_redutor}:

\begin{align*}
  F_{r_t} & = \frac{\num{6611.85} + \num{978.48}}{2} + \num{2537.78} \cdot 10 \\
  F_{r_t} & = \boxed{\SI{29172.97}{\newton}}
\end{align*}

Verificando a força radial de saída admitida pelo redutor adotado:

\begin{align*}
  F_r                   & > F_{r_t}                                        \\
  \SI{88}{\kilo\newton} & > \SI{29.17}{\kilo\newton} \rightarrow \text{OK}
\end{align*}
